\documentclass[11pt,a4paper]{article}
\usepackage{dsmorgan}

\title{The Golden Ratio is Irrational}
\subtitle{Three proofs: descent, rational roots and continued fractions}
\author{Drew M. Schaffner}
\date{Draft, \today}

\begin{document}
\maketitle

\section{Statement}

\begin{theorem}[Irrationality of $\varphi$]\label{thm:main}
  The golden ratio $\varphi = \dfrac{1 + \sqrt5}{2}$ is irrational.
\end{theorem}

\section{Setting}

The number $\varphi$ is the positive root of
\begin{equation}\label{eq:poly}
  x^2 - x - 1 = 0, \qquad\text{equivalently}\qquad \varphi = 1 + \frac{1}{\varphi}.
\end{equation}
Since $2 < \sqrt5 < 3$, we have $1 < \varphi < 2$.

\section{Idea of the proof}

The relation $\varphi = 1 + 1/\varphi$ is self-similar: removing a unit from $\varphi$
and inverting gives $\varphi$ back. A fraction cannot behave like this, because each
step strictly shrinks its denominator.

\section{First proof: descent}

\begin{proof}
  Suppose $\varphi$ is rational. Among all ways of writing $\varphi = p/q$ with
  integers $p, q \geq 1$, choose one with $q$ least, by well-ordering. By
  \eqref{eq:poly},
  \[
    \varphi = \frac{1}{\varphi - 1} = \frac{1}{p/q - 1} = \frac{q}{p - q}.
  \]
  Since $1 < p/q < 2$, we have $0 < p - q < q$. This writes $\varphi$ with the positive
  denominator $p - q < q$, contradicting the choice of $q$.
\end{proof}

\section{Second proof: rational roots}

\begin{lemma}[Rational root test]\label{lem:rrt}
  If $p/q$ is in lowest terms and is a root of $a x^2 + b x + c$ with integer
  coefficients, then $p \mid c$ and $q \mid a$.
\end{lemma}

\begin{proof}
  Multiplying $a(p/q)^2 + b(p/q) + c = 0$ by $q^2$ gives $ap^2 + bpq + cq^2 = 0$. So
  $p \mid cq^2$. Since $\gcd(p, q) = 1$, also $\gcd(p, q^2) = 1$, so B\'ezout gives
  $1 = px + q^2y$ and $c = pcx + cq^2y$ is divisible by $p$. Similarly, $q \mid ap^2$
  gives $q \mid a$.
\end{proof}

\begin{proof}[Second proof of \cref{thm:main}]
  By \cref{lem:rrt}, a rational root of $x^2 - x - 1$ has $p \mid 1$ and $q \mid 1$, so
  it is $\pm 1$. But $1 - 1 - 1 = -1$ and $1 + 1 - 1 = 1$, so neither is a root.
\end{proof}

\section{Third proof: continued fractions}

For $x > 0$, the \emph{continued fraction algorithm} sets $x_0 = x$, records
$a_j = \lfloor x_j \rfloor$, and continues with $x_{j+1} = 1/(x_j - a_j)$, stopping if
$x_j$ is an integer.

\begin{lemma}[Rationals terminate]\label{lem:terminate}
  If $x$ is rational, the algorithm stops after finitely many steps.
\end{lemma}

\begin{proof}
  If $x_j = p/q$ with $q \geq 1$ is not an integer, write $p = a_j q + r$ with
  $0 < r < q$. Then $x_{j+1} = q/r$ has denominator $r < q$. Denominators are positive
  integers that strictly decrease, so the process stops.
\end{proof}

\begin{proof}[Third proof of \cref{thm:main}]
  For $x_0 = \varphi$: $a_0 = 1$ because $1 < \varphi < 2$, and
  $x_1 = 1/(\varphi - 1) = \varphi$ by \eqref{eq:poly}. So $x_j = \varphi$ for every
  $j$, and the algorithm never stops:
  \[
    \varphi = 1 + \cfrac{1}{1 + \cfrac{1}{1 + \cfrac{1}{1 + \dotsb}}}.
  \]
  By \cref{lem:terminate}, $\varphi$ is irrational.
\end{proof}

\section{Remarks}

\begin{remark}
  The three proofs are one proof in different clothing. The step $p/q \mapsto q/(p-q)$
  in the first proof is one step of the continued fraction algorithm, and it is the
  Euclidean algorithm run on $(p, q)$.
\end{remark}

\begin{remark}
  Truncating the continued fraction gives $1, 2, \frac32, \frac53, \frac85, \dotsc$,
  the ratios $F_{j+1}/F_j$ of consecutive Fibonacci numbers. If $r_j = F_{j+1}/F_j$,
  then $r_{j+1} = 1 + 1/r_j$, the same recursion as \eqref{eq:poly}. One can show
  $\abs{r_j - \varphi} < 1/F_j^2$, so these fractions approximate $\varphi$, but only
  slowly: every partial quotient is $1$, the smallest possible.
\end{remark}

\begin{remark}
  This sharpens the treatment in the Foundations notes, which reduces the problem to
  $\sqrt5$ and residues mod $5$. Here no fact about $\sqrt5$ is needed beyond
  $2 < \sqrt5 < 3$.
\end{remark}

\end{document}
